% =====================================================================
%  Lecture 01 -- GSM, GPRS and EDGE                      (Persian edition)
%  Companion: lecture-01-en.tex.  Build: make lecture LECTURE=01
% =====================================================================
\documentclass[11pt,a4paper]{article}
\providecommand\ACNroot{../..}
\input{\ACNroot/templates/lecture-fa}

\ACNset{number}{01}
\ACNset{title}{جی‌اس‌ام، جی‌پی‌آر‌اس و اج}
\ACNset{shorttitle}{\lr{GSM}، \lr{GPRS} و \lr{EDGE}}
\ACNset{titleother}{GSM, GPRS and EDGE}
\ACNset{part}{بخش اول -- مبانی نسل دوم و سوم (درس‌های ۱ تا ۴)}
\ACNset{time}{۶ تا ۸ ساعت (۳ ساعت مطالعه، ۳ تا ۵ ساعت تمرین)}
\ACNset{level}{مقدماتی تا متوسط}
\ACNset{prereq}{دسی‌بل، مدولاسیون دیجیتال پایه، \lr{TDMA/FDMA}، مدل \lr{OSI}}
\ACNset{gen}{2}

\begin{document}
\ACNcover

% ---------------------------------------------------------------------
\section{اهداف یادگیری}\label{sec:objectives}
پس از مطالعهٔ این درس باید بتوانید:
\begin{objectives}
  \item معماری مرجع \acro{GSM} را \textbf{رسم کنید}، همهٔ عناصر شبکه
        (\lr{MS, BTS, BSC, MSC, VLR, GMSC, HLR, AuC, EIR}) را نام ببرید و واسط میان
        هر دو عنصر (\lr{Um, Abis, A, C, D, E, F}) را مشخص کنید.
  \item زمان‌بندی رادیویی \acro{GSM} (دورهٔ بیت، بُرست، قاب \acro{TDMA} و چندقاب ۲۶تایی) را
        از نرخ ناخالص \qty{270.833}{\kilo\bit\per\second} \textbf{استخراج کنید} و نرخ
        ناخالص \qty{22.8}{\kilo\bit\per\second} کانال ترافیک تمام‌نرخ را \textbf{محاسبه کنید}.
  \item احراز هویت چالش--پاسخ \acro{GSM} با سه‌تایی $(\text{RAND},\text{SRES},K_c)$ را
        \textbf{توضیح دهید} و ضعف‌های امنیتی آن را شناسایی کنید.
  \item لایهٔ افزودهٔ \acro{GPRS} (\lr{PCU, SGSN, GGSN, Gb, Gn, Gi, Gr}) و زمینهٔ \acro{PDP} را
        \textbf{توصیف کنید} و فعال‌سازی زمینهٔ \acro{PDP} را گام‌به‌گام \textbf{دنبال کنید}.
  \item گذردهی هر کاربر در \acro{GPRS} و \acro{EDGE} را از روی طرح کدگذاری و تعداد شیارهای
        زمانی \textbf{محاسبه کنید} و مصالحهٔ میان استحکام و نرخ را \textbf{ارزیابی کنید}.
  \item اندازهٔ خوشه، نمای افت مسیر و تداخل هم‌کانال را از طریق مدل $C/I$ ردیف اول
        \textbf{به هم ربط دهید}.
\end{objectives}

% ---------------------------------------------------------------------
\section{دانش پیش‌نیاز}\label{sec:prereq}
این درس فرض می‌کند با محاسبات دسی‌بلی مرور پیش‌نیازها
($x_{\mathrm{dB}} = 10\log_{10}x$)، مفهوم نماد مدولاسیون و پهنای باند، و تفاوت
دسترسی چندگانه با تقسیم فرکانس (\acro{FDMA}) و تقسیم زمان (\acro{TDMA}) آشنا هستید.
اگر در هر یک از این موارد تردید دارید، پیش از ادامه آن‌ها را مرور کنید؛ همهٔ نتایج
کمّی این درس بر پایهٔ همین مفاهیم ساخته می‌شوند.

\paragraph{واژه‌های کلیدی} جدول زیر واژگان به‌کاررفته در این درس را همراه با معادل
انگلیسی آن‌ها نشان می‌دهد تا دو نسخهٔ درس را بتوان در کنار هم خواند.
\begin{keyterms}
ایستگاه موبایل & \lr{Mobile Station (MS)} & گوشی؛ یعنی تجهیزات موبایل (\acro{ME}) به‌اضافهٔ کارت \acro{SIM}. \\
ایستگاه فرستنده‌گیرندهٔ پایه & \lr{BTS} & سایت رادیویی: فرستنده‌ها، گیرنده‌ها و آنتن‌های یک یا چند سلول. \\
کنترل‌گر ایستگاه پایه & \lr{BSC} & منابع رادیویی و تحویل‌های تعداد زیادی \acro{BTS} را مدیریت می‌کند. \\
مرکز سوئیچ موبایل & \lr{MSC} & سوئیچ مداری تماس‌ها؛ سیگنالینگ تحرک را نیز اجرا می‌کند. \\
ثبات مکان خانگی & \lr{HLR} & پایگاه دادهٔ دائمی مشترکان در اپراتور خانگی. \\
ثبات مکان مهمان & \lr{VLR} & نسخهٔ موقت داده‌های مشترکانی که در محدودهٔ یک \acro{MSC} هستند. \\
مرکز احراز هویت & \lr{AuC} & کلید محرمانهٔ $K_i$ را نگه می‌دارد و سه‌تایی‌ها را تولید می‌کند. \\
شیار زمانی & \lr{Timeslot (TS)} & یکی از هشت موقعیت \acro{TDMA} روی یک حامل. \\
بُرست & \lr{Burst} & شکل‌موج رادیویی ارسال‌شده در یک شیار زمانی. \\
زمینهٔ \acro{PDP} & \lr{PDP context} & وضعیت نشست بسته‌ای: نشانی، \acro{APN}، کیفیت خدمات و شناسه‌های تونل. \\
گره‌های پشتیبان \acro{GPRS} & \lr{SGSN / GGSN} & گره‌های خدمت‌دهنده و دروازهٔ هستهٔ بسته‌ای \acro{GPRS}. \\
طرح کدگذاری & \lr{CS / MCS} & ترکیب مدولاسیون و نرخ کد کانال. \\
\end{keyterms}

% ---------------------------------------------------------------------
\section{مقدمه و انگیزه}\label{sec:intro}
در دههٔ ۱۹۸۰ اروپای غربی چندین سامانهٔ سلولی آنالوگ ناسازگار با یکدیگر داشت. گوشی
خریداری‌شده در یک کشور در آن سوی مرز بی‌مصرف بود، هر سامانه صنعت تجهیزات خاص خود را
لازم داشت و صدای \acro{FM} آنالوگ هیچ محافظتی در برابر شنود نداشت. بنابراین مسئلهٔ مهندسی‌ای
که \acro{GSM}\footnote{در اصل \lr{\emph{Groupe Spécial Mobile}} و بعدها
\lr{\emph{Global System for Mobile communications}}.} برای حل آن طراحی شد تنها یک مسئلهٔ
رادیویی نبود؛ یک مسئلهٔ \emph{سامانه‌ای} بود:
\begin{itemize}
  \item \textbf{هم‌کنش‌پذیری.} واسط‌های باز و استاندارد به اپراتور اجازه می‌داد ایستگاه‌های
        پایه را از یک سازنده و سوئیچ‌ها را از سازندهٔ دیگری بخرد.
  \item \textbf{رومینگ.} هویت مشترک مستقل از گوشی (کارت \acro{SIM}) و پایگاه دادهٔ خانگی
        که از هر شبکهٔ میزبانی قابل دسترسی باشد.
  \item \textbf{بازدهی طیفی و ظرفیت.} کدگذاری دیجیتال گفتار، کدگذاری کانال و \acro{TDMA}
        اجازه می‌دهند چند کاربر یک حامل را به اشتراک بگذارند و تداخل بیشتری را تحمل کنند؛
        در نتیجه فرکانس‌ها را می‌توان فشرده‌تر بازاستفاده کرد.
  \item \textbf{امنیت.} احراز هویت مشترک و رمزنگاری پیوند رادیویی.
\end{itemize}
\acro{GSM} برای صدای سوئیچ مداری طراحی شد. وقتی تقاضای داده در دههٔ ۱۹۹۰ رشد کرد، همان
شبکهٔ رادیویی با یک لایهٔ افزودهٔ سوئیچ بسته‌ای گسترش یافت: \acro{GPRS}
(\lr{General Packet Radio Service}) و سپس \acro{EDGE}
(\lr{Enhanced Data rates for GSM Evolution})، که با مدولاسیون مرتبهٔ بالاتر نرخ بیت هر
شیار زمانی را افزایش داد. بررسی این سه گام در یک درس الگویی را نشان می‌دهد که در سراسر
این درس تکرار می‌شود: \textbf{هر نسل هر چه را بتواند دوباره به کار می‌گیرد و فقط
حداقل سازوکار تازه‌ای را می‌افزاید که مسئلهٔ جدید لازم دارد.}

% ---------------------------------------------------------------------
\section{مفاهیم بنیادی}\label{sec:fundamentals}

\subsection{سلول‌ها، بازاستفادهٔ فرکانس و دسترسی چندگانه}
شبکهٔ موبایل ناحیهٔ خدمت خود را به \emph{سلول‌ها} تقسیم می‌کند که هر یک را یک مجموعه
آنتن پوشش می‌دهد. چون طیف کمیاب است، فرکانس‌های حامل یکسان در سلول‌هایی که به اندازهٔ
کافی از هم دورند \emph{بازاستفاده} می‌شوند. \acro{GSM} دو روش دسترسی چندگانه را ترکیب می‌کند:
\begin{itemize}
  \item \textbf{\acro{FDMA}:} باند به حامل‌هایی با فاصلهٔ \qty{200}{\kilo\hertz} تقسیم می‌شود.
        در باند اصلی \lr{GSM-900}، فراسو بازهٔ \qtyrange{890}{915}{\mega\hertz} و فروسو بازهٔ
        \qtyrange{935}{960}{\mega\hertz} را اشغال می‌کند؛ نتیجه ۱۲۴ جفت حامل قابل استفاده با
        فاصلهٔ دوطرفهٔ \qty{45}{\mega\hertz} است \cite{ts45002}.
  \item \textbf{\acro{TDMA}:} هر حامل در زمان به قاب‌هایی با هشت شیار زمانی تقسیم می‌شود؛
        یک کاربر صوتی تمام‌نرخ در هر قاب یک شیار زمانی را اشغال می‌کند.
\end{itemize}

\begin{definitionbox}{کانال فیزیکی و کانال منطقی}
\emph{کانال فیزیکی} یک شمارهٔ شیار زمانی روی یک حامل (یا روی دنباله‌ای از حامل‌ها در
پرش فرکانسی) است. \emph{کانال منطقی} نوعی از اطلاعات است --- ترافیک (\acro{TCH}) یا
کنترل (\lr{BCCH, CCCH, SDCCH, SACCH, FACCH}) --- که بر اساس یک ساختار چندقابی ثابت روی
کانال فیزیکی نگاشت می‌شود.
\end{definitionbox}

\subsection{جدا کردن هویت از تجهیزات}
\acro{GSM} \emph{مشترک} را از \emph{دستگاه} جدا می‌کند. مشترک با \acro{IMSI}
(\lr{International Mobile Subscriber Identity}: کد کشور، کد شبکه و شمارهٔ مشترک، حداکثر
۱۵ رقم) شناخته می‌شود که همراه با کلید محرمانهٔ ۱۲۸ بیتی $K_i$ روی \acro{SIM} ذخیره
شده است. دستگاه با \acro{IMEI} شناخته می‌شود. برای اینکه \acro{IMSI} روی هوا ارسال نشود،
\acro{VLR} یک هویت موقت به نام \acro{TMSI} اختصاص می‌دهد که هر جا ممکن باشد به‌جای آن
به کار می‌رود \cite{ts23003}.

\subsection{سوئیچ مداری در برابر سوئیچ بسته‌ای}
صدا در \acro{GSM} \emph{سوئیچ مداری} است: با برقراری تماس، یک شیار زمانی روی رادیو، یک
کانال روی پیوندهای \acro{Abis} و \acro{A} و یک مسیر درون \acro{MSC} برای کل مدت تماس
رزرو می‌شود، چه کسی صحبت کند و چه نکند. داده در \acro{GPRS} \emph{سوئیچ بسته‌ای} است:
بلوک‌های رادیویی تنها هنگامی تخصیص می‌یابند که کاربر داده‌ای برای ارسال داشته باشد و
کاربران زیادی شیارهای زمانی یکسانی را به اشتراک می‌گذارند. همین تفاوت دلیل آن است که
\acro{GPRS} به گره‌های جدید نیاز دارد و صرفاً نرم‌افزار جدیدی در \acro{MSC} کافی نیست؛
این موضوع در درس ۲ به‌تفصیل تحلیل می‌شود.

% ---------------------------------------------------------------------
\section{توضیحات فنی تفصیلی}\label{sec:technical}

\subsection{عناصر شبکهٔ \acro{GSM}}
\paragraph{ایستگاه موبایل (\acro{MS})} ایستگاه موبایل همان تجهیزات موبایل به‌اضافهٔ \acro{SIM}
است. \acro{SIM} الگوریتم احراز هویت \acro{A3} و الگوریتم تولید کلید \acro{A8} را اجرا می‌کند
تا $K_i$ هرگز از کارت خارج نشود.

\paragraph{ایستگاه فرستنده‌گیرندهٔ پایه (\acro{BTS})} \acro{BTS} فرستنده‌گیرنده‌های رادیویی
(\acro{TRX}) یک یا چند سلول را در بر دارد. مدولاسیون (\acro{GMSK})، کدگذاری کانال،
درهم‌آمیزی (\lr{interleaving})، رمزنگاری با \acro{A5} و اندازه‌گیری‌های رادیویی که مبنای
تصمیم‌های تحویل هستند در آن انجام می‌شوند. این عنصر عمداً ساده است: تصمیم‌ها را \acro{BSC} می‌گیرد.

\paragraph{کنترل‌گر ایستگاه پایه (\acro{BSC})} یک \acro{BSC} ده‌ها تا صدها \acro{BTS} را کنترل
می‌کند. کانال‌های رادیویی را تخصیص می‌دهد، کنترل توان را مدیریت می‌کند و بیشتر تحویل‌ها
میان سلول‌های خود را تصمیم می‌گیرد. مجموعهٔ \acro{BTS}ها و \acro{BSC} آن‌ها \emph{زیرسامانهٔ
ایستگاه پایه} (\acro{BSS}) را تشکیل می‌دهد. تبدیل کدک گفتار میان کدک رادیویی
\qty{13}{\kilo\bit\per\second} و قالب \acro{PCM} \qty{64}{\kilo\bit\per\second} هسته، در
واحد تبدیل کدک (\acro{TRAU}) انجام می‌شود که برای صرفه‌جویی در ظرفیت انتقال معمولاً در
محل \acro{MSC} قرار می‌گیرد.

\paragraph{مرکز سوئیچ موبایل (\acro{MSC}) و \acro{VLR}} \acro{MSC} یک مرکز تلفن است که
کارکردهای تحرک به آن افزوده شده‌اند: تماس‌ها را برقرار و آزاد می‌کند، تحویل میان
\acro{BSC}ها را کنترل می‌کند و به‌روزرسانی مکان و احراز هویت را اجرا می‌کند. \acro{VLR}
(که در عمل درون \acro{MSC} قرار دارد) داده‌های اشتراک و سه‌تایی‌های امنیتی همهٔ مشترکانی
را که اکنون در محدودهٔ آن هستند نگه می‌دارد.

\paragraph{\acro{MSC} دروازه (\acro{GMSC})} تماس‌هایی که از شبکهٔ ثابت
(\lr{PSTN/ISDN}) می‌رسند به یک \acro{GMSC} هدایت می‌شوند که از \acro{HLR} می‌پرسد مشترک
مقصد کجاست.

\paragraph{\acro{HLR}، \acro{AuC} و \acro{EIR}} \acro{HLR} نمایهٔ دائمی اشتراک و نشانی
\acro{VLR}ای را که اکنون به مشترک خدمت می‌دهد ذخیره می‌کند. \acro{AuC} کلید $K_i$ را نگه
می‌دارد و سه‌تایی‌های احراز هویت را محاسبه می‌کند. \acro{EIR} فهرست‌های سفید، خاکستری و
سیاه \acro{IMEI}ها را نگه می‌دارد.

\subsection{واسط‌ها}
\begin{table}[htbp]
\centering\small
\caption{واسط‌های اصلی \acro{GSM} و پروتکل‌های آن‌ها \cite{ts23002}.}\label{tab:interfaces}
\rowcolors{2}{acnmist}{white}
\begin{tabularx}{\linewidth}{@{}l l l Y@{}}
\toprule
\rowcolor{white}واسط & میان & پروتکل اصلی & نقش \\
\midrule
\lr{Um}   & \lr{MS -- BTS}   & \lr{LAPDm, RR, MM, CC} & واسط رادیویی، \acro{TDMA/FDMA}، \acro{GMSK} \\
\lr{Abis} & \lr{BTS -- BSC}  & \lr{LAPD over E1/T1}   & انتقال ترافیک و کنترل \acro{BTS} \\
\lr{A}    & \lr{BSC -- MSC}  & \lr{BSSAP over SS7}    & سیگنالینگ تماس و تحرک، مدارهای \qty{64}{\kilo\bit\per\second} \\
\lr{B}    & \lr{MSC -- VLR}  & \lr{MAP}               & در عمل داخلی (\acro{MSC/VLR} یکپارچه) \\
\lr{C}    & \lr{GMSC -- HLR} & \lr{MAP}               & پرسش مسیریابی برای تماس‌های ورودی \\
\lr{D}    & \lr{VLR -- HLR}  & \lr{MAP}               & به‌روزرسانی مکان، داده‌های مشترک \\
\lr{E}    & \lr{MSC -- MSC}  & \lr{MAP + ISUP}        & تحویل میان \acro{MSC}ها، مسیریابی تماس \\
\lr{F}    & \lr{MSC -- EIR}  & \lr{MAP}               & بررسی \acro{IMEI} \\
\bottomrule
\end{tabularx}
\end{table}

ارزش یک واسط باز همان‌قدر که فنی است اقتصادی نیز هست: چون واسط \acro{A} به‌طور کامل
استاندارد شده است، اپراتور می‌تواند \acro{BSS} هر سازنده‌ای را به \acro{MSC} هر سازندهٔ
دیگری وصل کند. واسط \acro{Abis} نیز استاندارد شد، اما در عمل اغلب افزونه‌های اختصاصی
سازنده برای بهره‌برداری و نگهداری دارد.

\subsection{واسط رادیویی \acro{Um}}
\acro{GSM} از مدولاسیون کلیدزنی کمینهٔ جابه‌جایی گاوسی (\acro{GMSK}) با حاصل‌ضرب پهنای
باند در زمان $BT = 0.3$ استفاده می‌کند. \acro{GMSK} مدولاسیونی با پوش ثابت است؛ بنابراین
تقویت‌کنندهٔ توان گوشی می‌تواند نزدیک اشباع کار کند و این کار باتری را صرفه‌جویی می‌کند.
در هر نماد یک بیت حمل می‌شود، با نرخ ناخالص
$13/48\,\mathrm{MHz} \approx \qty{270.833}{\kilo\bit\per\second}$. ساختار زمانی‌ای که از این
نرخ به دست می‌آید در \cref{fig:tdma} رسم و در \cref{sec:models} محاسبه شده است.

یک بُرست عادی دو بلوک ۵۷ بیتی داده، یک دنبالهٔ آموزشی ۲۶ بیتی در وسط (گیرنده از آن برای
تخمین کانال چندمسیره در همسان‌ساز استفاده می‌کند)، دو پرچم ربایش، بیت‌های انتهایی و یک
دورهٔ محافظ به طول ۸٫۲۵ دورهٔ بیت دارد که عدم قطعیت زمانی و شیب افزایش و کاهش توان را جذب
می‌کند. پرش فرکانسی (تا ۲۱۷ پرش در ثانیه، یک پرش در هر قاب \acro{TDMA}) محوشدگی و تداخل
را میان حامل‌ها میانگین می‌گیرد.

\subsection{تحرک: نواحی مکانی و \acro{TMSI}}
سلول‌ها در \emph{نواحی مکانی} (\acro{LA}) گروه‌بندی می‌شوند. \acro{MS} در حالت بیکار به
کانال پخش گوش می‌دهد و تنها هنگام ورود به یک \acro{LA} جدید، یا به‌صورت دوره‌ای،
\emph{به‌روزرسانی مکان} انجام می‌دهد. وقتی تماسی می‌رسد، شبکه \acro{MS} را در همهٔ
سلول‌های \acro{LA} فعلی آن فراخوانی (\lr{paging}) می‌کند. اندازهٔ \acro{LA} مصالحه‌ای است
میان سیگنالینگ به‌روزرسانی (\acro{LA}های کوچک) و بار فراخوانی (\acro{LA}های بزرگ)؛ این
مصالحه در درس‌های ۴ و ۱۲ به‌صورت کمّی بررسی می‌شود.

\subsection{امنیت: احراز هویت و رمزنگاری}
احراز هویت در \acro{GSM} یک پروتکل \emph{چالش--پاسخ} است \cite{ts43020}. \acro{AuC} یک عدد
تصادفی ۱۲۸ بیتی \lr{RAND} انتخاب و مقادیر زیر را محاسبه می‌کند
\begin{subequations}\label{eq:a3a8}
\begin{align}
  \text{SRES} &= \mathrm{A3}(K_i,\text{RAND}), \qquad &&\text{(32 bits)} \label{eq:a3}\\
  K_c &= \mathrm{A8}(K_i,\text{RAND}), \qquad &&\text{(64 bits)} \label{eq:a8}
\end{align}
\end{subequations}
و سه‌تایی $(\text{RAND},\text{SRES},K_c)$ را برای \acro{VLR} می‌فرستد. \acro{VLR} تنها
\lr{RAND} را به \acro{MS} می‌فرستد؛ \acro{SIM} محاسبهٔ \cref{eq:a3} را تکرار و نتیجه را
بازمی‌گرداند. اگر نتیجه با \lr{SRES} برابر باشد مشترک معتبر است و اکنون هر دو طرف $K_c$ را
به اشتراک دارند که کلید رمز جریانی \acro{A5} روی پیوند رادیویی است. توجه کنید این طراحی
چه چیزی را تأمین می‌کند و چه چیزی را نه:
\begin{itemize}
  \item $K_i$ هرگز از هیچ واسطی عبور نمی‌کند و \lr{RAND} بدون $K_i$ بی‌فایده است.
  \item احراز هویت \emph{یک‌طرفه} است: شبکه برای \acro{MS} احراز هویت نمی‌شود؛ بنابراین
        یک ایستگاه پایهٔ جعلی می‌تواند خود را به جای شبکه جا بزند.
  \item رمزنگاری در \acro{BTS} پایان می‌یابد؛ پیوندهای \acro{Abis}، \acro{A} و هسته با
        \acro{A5} محافظت نمی‌شوند.
  \item $K_c$ تنها ۶۴ بیت دارد و برای رمزهای \acro{A5/1} و \acro{A5/2} حمله‌های رمزشکنی
        منتشر شده است؛ \acro{A5/2} کنار گذاشته شد و \acro{A5/3} (مبتنی بر \acro{KASUMI})
        بعدها افزوده شد.
\end{itemize}
\acro{UMTS} (درس ۳) و \acro{LTE} (درس‌های ۹ و ۱۱) این ضعف‌ها را با احراز هویت دوطرفه،
حفاظت از یکپارچگی و کلیدهای بلندتر برطرف می‌کنند.

\subsection{\acro{GPRS}: یک لایهٔ افزودهٔ بسته‌ای}
\acro{GPRS} سوئیچ بسته‌ای را می‌افزاید و در عین حال حامل‌های رادیویی، سایت‌ها و پایگاه‌های
دادهٔ مشترک \acro{GSM} را دوباره به کار می‌گیرد \cite{ts23060}. سه کارکرد جدید لازم است:
\begin{itemize}
  \item \textbf{واحد کنترل بسته (\acro{PCU})} در \acro{BSS} که لایهٔ بسته‌محور
        \acro{RLC/MAC} را پیاده می‌کند و بلوک‌های رادیویی را میان کاربران زمان‌بندی می‌کند.
  \item \textbf{گره پشتیبان خدمت‌دهندهٔ \acro{GPRS} (\acro{SGSN})}، همتای بسته‌ای \acro{MSC}:
        مدیریت تحرک (نواحی مسیریابی)، احراز هویت، رمزنگاری (در \acro{GPRS} رمزنگاری میان
        \acro{MS} و \acro{SGSN} در لایهٔ \acro{LLC} انجام می‌شود) و صورت‌حساب.
  \item \textbf{گره پشتیبان دروازهٔ \acro{GPRS} (\acro{GGSN})}، مسیریاب مرزی به سوی
        شبکه‌های داده‌ای بیرونی: نشانی \acro{IP} کاربر را تخصیص می‌دهد و لنگرگاه نشست است.
\end{itemize}
\acro{SGSN} از طریق \iface{Gb} (در اصل \lr{Frame Relay}) به \acro{BSS}، از طریق \iface{Gn}
و با پروتکل تونل‌زنی \acro{GPRS} (\acro{GTP}) به \acro{GGSN}، و از طریق \iface{Gr} و با
\acro{MAP} به \acro{HLR} وصل می‌شود. \acro{GGSN} از طریق \iface{Gi} به اینترنت یا شبکهٔ
سازمانی متصل است. ایدهٔ \emph{تونل} میان یک گره خدمت‌دهنده و یک گره دروازه ماندگارترین
مفهومی است که \acro{GPRS} معرفی کرد: این مفهوم به‌صورت واسط \lr{S5/S8} در \acro{LTE} و
واسط \lr{N9} در \lr{5G} باقی مانده است.

در سمت رادیو، \acro{GPRS} کانال‌های دادهٔ بسته‌ای (\acro{PDCH}) را تعریف می‌کند. \acro{PDCH}
یک شیار زمانی معمولی است که به ترافیک بسته‌ای اختصاص یافته است؛ داده در \emph{بلوک‌های
رادیویی} شامل چهار بُرست متوالی ارسال می‌شود (به‌طور میانگین یک بلوک در هر
\qty{20}{\milli\second}). چهار طرح کدگذاری میان حفاظت و نرخ مصالحه برقرار می‌کنند:
\lr{CS-1} (\qty{9.05}{\kilo\bit\per\second})، \lr{CS-2} (13.4)، \lr{CS-3} (15.6) و \lr{CS-4}
(\qty{21.4}{\kilo\bit\per\second}) برای هر شیار زمانی. یک \acro{MS} چندشیاره می‌تواند
هم‌زمان از چند شیار زمانی استفاده کند.

\subsection{حالت‌های مدیریت تحرک در \acro{GPRS}}
\acro{MS} در \acro{GPRS} در یکی از سه حالت مدیریت تحرک است:
\textbf{\lr{IDLE}} (متصل نیست؛ شبکه \acro{MS} را نمی‌شناسد)،
\textbf{\lr{STANDBY}} (متصل است؛ مکان در سطح ناحیهٔ مسیریابی معلوم است؛ پیش از تحویل
داده باید فراخوانی شود) و \textbf{\lr{READY}} (مکان در سطح سلول معلوم است؛ داده بی‌درنگ
قابل ارسال است). یک زمان‌سنج پس از مدتی بی‌فعالیتی \acro{MS} را از \lr{READY} به
\lr{STANDBY} بازمی‌گرداند. همین ایده --- دستگاه را در حالت بیکار با هزینهٔ کم دسترس‌پذیر
نگه دار و تنها هنگام جریان ترافیک سیگنالینگ خرج کن --- به‌صورت \lr{ECM-IDLE/CONNECTED}
در \acro{LTE} و \lr{CM-IDLE/CONNECTED} در \lr{5G} دوباره ظاهر می‌شود.

\subsection{\acro{EDGE}: مدولاسیون مرتبهٔ بالاتر روی همان حامل}
\acro{EDGE} نرخ نماد \acro{GSM} (\lr{270.833\,ksymbol/s})، حامل \qty{200}{\kilo\hertz} و ساختار
\acro{TDMA} را حفظ می‌کند، اما \lr{8-PSK} را معرفی می‌کند که به‌جای یک بیت، سه بیت در هر
نماد حمل می‌کند. \acro{GPRS} بهبودیافته (\acro{EGPRS}) نُه طرح مدولاسیون و کدگذاری تعریف
می‌کند: \lr{MCS-1} تا \lr{MCS-4} از \acro{GMSK} و \lr{MCS-5} تا \lr{MCS-9} از \lr{8-PSK}
استفاده می‌کنند و نرخ هر شیار زمانی از \qty{8.8}{\kilo\bit\per\second} تا
\qty{59.2}{\kilo\bit\per\second} است. دو سازوکار نرخ‌های بالاتر را قابل استفاده می‌کنند:
\begin{itemize}
  \item \textbf{تطبیق پیوند:} شبکه \acro{MCS} را بر اساس اندازه‌گیری‌های کیفیت کانال
        انتخاب می‌کند و وقتی کاربر در لبهٔ سلول است به طرحی مقاوم‌تر می‌رود.
  \item \textbf{افزونگی افزایشی:} اگر بلوکی با خطا دریافت شود، ارسال مجدد سوراخ‌کاری
        (\lr{puncturing}) متفاوتی از همان بیت‌های کدشده را حمل می‌کند و گیرنده هر دو نسخه را
        ترکیب می‌کند. این شکل اولیه‌ای از \acro{HARQ} است که در \acro{HSPA}، \acro{LTE} و
        \acro{NR} به کار می‌رود.
\end{itemize}

% ---------------------------------------------------------------------
\section{مدل‌های ریاضی}\label{sec:models}

\subsection{زمان‌بندی رادیویی \acro{GSM}}\label{sec:timing}
\paragraph{مسئلهٔ مهندسی} همهٔ پارامترهای زمانی \acro{GSM} --- طول بُرست، مدت قاب،
بیشینهٔ شعاع سلول --- از یک عدد، یعنی نرخ بیت ناخالص، نتیجه می‌شوند. آن‌ها را استخراج
می‌کنیم تا ببینیم مقادیر آشنا از کجا می‌آیند.

\paragraph{مدل سامانه و فرض‌ها} یک حامل، \acro{GMSK} با یک بیت در هر نماد، نرخ ناخالص
$R = f_{\mathrm{ref}}/48$ با ساعت مرجع $f_{\mathrm{ref}} = \qty{13}{\mega\hertz}$. هر شیار
زمانی ۱۵۶٫۲۵ دورهٔ بیت طول می‌کشد و هر قاب \acro{TDMA} هشت شیار زمانی دارد.

\paragraph{استخراج} دورهٔ بیت و مدت‌های حاصل از آن عبارت‌اند از
\begin{align}
  T_{\mathrm{b}} &= \frac{1}{R} = \frac{48}{13\times10^{6}}\ \mathrm{s}
     \approx \qty{3.692}{\micro\second}, \label{eq:tb}\\
  T_{\mathrm{slot}} &= 156.25\,T_{\mathrm{b}} = \frac{15}{26}\ \mathrm{ms}
     \approx \qty{576.9}{\micro\second}, \label{eq:tslot}\\
  T_{\mathrm{frame}} &= 8\,T_{\mathrm{slot}} = \frac{60}{13}\ \mathrm{ms}
     \approx \qty{4.615}{\milli\second}, \label{eq:tframe}\\
  T_{26} &= 26\,T_{\mathrm{frame}} = \qty{120}{\milli\second}. \label{eq:t26}
\end{align}
\paragraph{نرخ ناخالص کانال ترافیک تمام‌نرخ} در هر چندقاب ۲۶تایی، یک \acro{TCH} تمام‌نرخ ۲۴
قاب را برای ترافیک، یک قاب را برای \acro{SACCH} به کار می‌برد و یک قاب را بیکار می‌گذارد.
هر بُرست ترافیک $2\times 57 = 114$ بیت کدشده حمل می‌کند؛ بنابراین
\begin{equation}\label{eq:tchrate}
  R_{\mathrm{TCH/F}} = \frac{114\ \mathrm{bit}\times 24}{T_{26}}
  = \frac{2736\ \mathrm{bit}}{0.120\ \mathrm{s}} = \qty{22.8}{\kilo\bit\per\second}.
\end{equation}
خروجی \qty{13}{\kilo\bit\per\second} کدک گفتار با کدگذاری کانالی محافظت می‌شود که آن را
دقیقاً به همین \qty{22.8}{\kilo\bit\per\second} گسترش می‌دهد.

\paragraph{تفسیر فیزیکی} نسبت $R_{\mathrm{TCH/F}}/R = 22.8/270.833 \approx 8.4\%$ نزدیک به
$1/8$ است منهای سهمی از قاب که صرف دنبالهٔ آموزشی، بیت‌های انتهایی، زمان محافظ و سیگنالینگ
می‌شود. بنابراین \acro{GSM} حدود یک‌چهارم ظرفیت ناخالص خود را صرف سربار می‌کند --- عددی
که هنگام مقایسه با سامانه‌های \acro{OFDM} در درس‌های ۶ و ۱۴ باید به خاطر داشت.

\paragraph{بازدهی طیفی} با $B = \qty{200}{\kilo\hertz}$ بازدهی طیفی ناخالص
$\eta = R/B \approx \qty{1.35}{bit/s/Hz}$ است. کران شانون \cite{shannon1948mathematical}
\begin{equation}\label{eq:shannon}
  C = B\log_2\!\left(1+\SNR\right)
\end{equation}
که در آن $C$ ظرفیت بر حسب \lr{bit/s}، $B$ پهنای باند بر حسب \lr{Hz} و $\SNR$ نسبت خطی
سیگنال به نویز است، نشان می‌دهد که $1.35$ بیت بر ثانیه بر هرتز در اصل می‌توانست در
$\SNR = 2^{1.354}-1 \approx 1.56$ (حدود \qty{1.9}{\dB}) حمل شود. \acro{GSM} در \acro{SNR}
بسیار بالاتری کار می‌کند (حدود \qty{9}{\dB} نسبت حامل به تداخل)، زیرا به‌جای کدهای
نزدیک به ظرفیت از مدولاسیون و کدگذاری ساده و مقاوم استفاده می‌کند. \Cref{eq:shannon} یک
کانال ایدهٔ‌ال با پهنای باند محدود و نویز سفید گاوسی جمع‌شونده (\acro{AWGN}) و پیچیدگی
کدگذاری نامحدود را فرض می‌کند؛ این یک کران بالاست، نه پیش‌بینی گذردهی.

\subsection{پیش‌فرستی زمانی و بیشینهٔ شعاع سلول}
\paragraph{مسئله} \acro{BTS} انتظار دارد هر بُرست درون شیار زمانی خودش برسد. یک \acro{MS}
دوردست فروسو را دیرتر می‌بیند و فراسوی آن نیز دوباره با تأخیر می‌رسد؛ بنابراین باید به
اندازهٔ تأخیر رفت‌وبرگشت زودتر ارسال کند.
\paragraph{مدل} انتشار در خط دید با سرعت $c$؛ پیش‌فرستی زمانی (\acro{TA}) با گام‌هایی به
اندازهٔ یک دورهٔ بیت اعلام می‌شود، $\mathrm{TA}\in\{0,1,\dots,63\}$.
\paragraph{استخراج} تأخیر رفت‌وبرگشت $2d/c$ باید با $\mathrm{TA}\cdot T_{\mathrm{b}}$ جبران
شود؛ پس هر گام \acro{TA} متناظر با فاصلهٔ زیر است
\begin{equation}\label{eq:tastep}
  \Delta d = \frac{c\,T_{\mathrm{b}}}{2} = \frac{\qty{2.998e8}{\metre\per\second}\times\qty{3.692}{\micro\second}}{2}
  \approx \qty{553}{\metre},
  \qquad d_{\max} = 63\,\Delta d \approx \qty{34.9}{\kilo\metre}.
\end{equation}
\paragraph{تفسیر و محدودیت‌ها} حد آشنای «۳۵ کیلومتر» سلول \acro{GSM} پیامد میدان ۶ بیتی
\acro{TA} است، نه بودجهٔ پیوند. سلول‌های با برد گسترده این حد را با استفاده از دو شیار
زمانی مجاور برای هر کاربر، به بهای کاهش ظرفیت، دور می‌زنند. همین استدلال با گامی ریزتر،
پیش‌فرستی زمانی \acro{LTE} را در درس ۸ به دست می‌دهد.

\subsection{تداخل هم‌کانال و اندازهٔ خوشه}\label{sec:reuse}
\paragraph{مسئله} حامل‌ها را تا چه حد می‌توان فشرده بازاستفاده کرد؟ فاصلهٔ بازاستفادهٔ
کوچک‌تر یعنی حامل بیشتر در هر سلول، اما تداخل قوی‌تر.
\paragraph{مدل و فرض‌ها} سلول‌های شش‌ضلعی با شعاع $R_{\mathrm{c}}$، خوشه‌هایی با
$N = i^2 + ij + j^2$ سلول ($i$ و $j$ اعداد صحیح نامنفی)، افت مسیر توانی
$\Prx \propto d^{-n}$ با نمای $n$، کاربری در لبهٔ سلول، و تنها شش سلول هم‌کانال ردیف اول
در فاصلهٔ بازاستفادهٔ $D = R_{\mathrm{c}}\sqrt{3N}$ \cite{rappaport2002wireless}.
\paragraph{استخراج} سیگنال مطلوب $C \propto R_{\mathrm{c}}^{-n}$ است و هر یک از شش تداخل‌کننده
تقریباً $D^{-n}$ سهم دارد؛ بنابراین
\begin{equation}\label{eq:ci}
  \frac{C}{I} \approx \frac{R_{\mathrm{c}}^{-n}}{6\,D^{-n}}
  = \frac{1}{6}\left(\frac{D}{R_{\mathrm{c}}}\right)^{n}
  = \frac{\left(\sqrt{3N}\right)^{n}}{6}.
\end{equation}
\paragraph{تفسیر} \Cref{eq:ci} با افزایش $N$ و $n$ بزرگ می‌شود: خوشهٔ بزرگ تداخل‌کننده‌ها
را دورتر می‌برد و افت مسیر تند آن‌ها را سریع‌تر تضعیف می‌کند. چون \acro{GSM} نسبت $C/I$ حدود
\qty{9}{\dB} را تحمل می‌کند، \cref{fig:reuse} نشان می‌دهد که $N=4$ (برای $n=4$ برابر
$13.8$\,dB) با حاشیهٔ مناسب کار می‌کند و $N=3$ برای $n=4$ ممکن ولی برای $n=3.5$ مرزی است.
در عمل اپراتورها بازاستفادهٔ فشرده را با پرش فرکانسی و کنترل توان ترکیب می‌کردند که تداخل
را میانگین‌گیری و کاهش می‌دهند.
\paragraph{محدودیت‌ها} این مدل محوشدگی سایه‌ای (که $C/I$ را چند دسی‌بل تغییر می‌دهد)،
بخش‌بندی سلول (که برخی تداخل‌کننده‌ها را حذف می‌کند) و این واقعیت را که همهٔ تداخل‌کننده‌ها
هم‌زمان فعال نیستند نادیده می‌گیرد. این یک تخمین اولیهٔ طراحی است، نه تضمین.

\begin{acnfigure}
  \acnfig{lecture-01/reuse-ci}
  \caption{نسبت $C/I$ ردیف اول از \cref{eq:ci} بر حسب اندازهٔ خوشه برای دو نمای افت مسیر.
  خط افقی نسبت حفاظتی تقریبی \acro{GSM} را نشان می‌دهد.}
  \label{fig:reuse}
\end{acnfigure}

\subsection{گذردهی بسته‌ای با چند شیار زمانی}
برای کاربری که $n_{\mathrm{TS}}$ شیار زمانی با نرخ $r_{\mathrm{CS}}$ برای هر شیار گرفته است،
بیشینهٔ نرخ دادهٔ لایهٔ کنترل پیوند رادیویی (\acro{RLC}) برابر است با
\begin{equation}\label{eq:gprsrate}
  R_{\mathrm{user}} = n_{\mathrm{TS}}\; r_{\mathrm{CS}}\;(1-p_{\mathrm{BLER}}),
\end{equation}
که در آن $p_{\mathrm{BLER}}$ نسبت خطای بلوک است و ضریب $(1-p_{\mathrm{BLER}})$ فرض می‌کند هر بلوک
خطادار یک بار و با موفقیت دوباره ارسال می‌شود (یک مدل ساده از \acro{ARQ}). طرح‌های کدگذاری
بالاتر $r_{\mathrm{CS}}$ را افزایش می‌دهند، اما در یک $C/I$ معین $p_{\mathrm{BLER}}$ را نیز بالا
می‌برند؛ بهترین طرح طرحی است که حاصل‌ضرب را بیشینه کند. \Cref{fig:rates} مقادیر بیشینه را
نشان می‌دهد.

\begin{acnfigure}
  \acnfig{lecture-01/gprs-rates}
  \caption{بیشینهٔ نرخ دادهٔ \acro{RLC} هر کاربر بر حسب تعداد شیارهای زمانی
  ($p_{\mathrm{BLER}}=0$) برای دو طرح \acro{GPRS} و دو طرح \acro{EGPRS}. گذردهی واقعی به دلیل
  ارسال مجدد، اشتراک میان کاربران و سربار لایه‌های بالاتر کمتر است.}
  \label{fig:rates}
\end{acnfigure}

% ---------------------------------------------------------------------
\section{معماری شبکه و نمودارها}\label{sec:architecture}
\Cref{fig:arch} شبکهٔ سوئیچ مداری \acro{GSM} و لایهٔ افزودهٔ بسته‌ای \acro{GPRS} را با هم
نشان می‌دهد. دو نکته شایان توجه است. نخست، \acro{BSS} مشترک است: همان \acro{BTS}
شیارهای صوتی را به سوی \acro{MSC} و شیارهای بسته‌ای را به سوی \acro{SGSN} می‌برد و
\acro{BSC} آن‌ها را از هم جدا می‌کند. دوم، \acro{HLR} مشترک است: یک رکورد اشتراک به هر دو
حوزه خدمت می‌دهد و به همین دلیل \acro{GPRS} بدون نیاز به تعریف دوبارهٔ مشترکان راه‌اندازی شد.

\begin{acnfigure}
  \acnfig{lecture-01/gsm-architecture}
  \caption{شبکهٔ \acro{GSM} همراه با لایهٔ افزودهٔ \acro{GPRS}. خاکستری توپر: پیوندهای انتقال؛
  فیروزه‌ای: پیوندهای سوئیچ بسته‌ای؛ آبی خط‌چین: سیگنالینگ \acro{MAP} به سوی \acro{HLR} و \acro{EIR}.}
  \label{fig:arch}
\end{acnfigure}

\begin{acnfigure}
  \resizebox{\linewidth}{!}{\acnfig{lecture-01/tdma-hierarchy}}
  \caption{ساختار زمانی \acro{GSM}. یک چندقاب ۲۶تایی (\qty{120}{\milli\second}) شامل ۲۴ قاب
  ترافیک، یک قاب \acro{SACCH} و یک قاب بیکار است؛ هر قاب \acro{TDMA} هشت شیار زمانی دارد؛
  یک بُرست عادی ۱۱۴ بیت داده حمل می‌کند.}
  \label{fig:tdma}
\end{acnfigure}

% ---------------------------------------------------------------------
\section{رویه‌های گام‌به‌گام}\label{sec:procedures}

\begin{procedurebox}{احراز هویت و آغاز رمزنگاری}\label{proc:auth}
\begin{steps}
  \item \acro{MS} یک \msg{Location Updating Request} (یا درخواست برقراری تماس) می‌فرستد و
        خود را با \acro{TMSI}، یا اگر \acro{TMSI} معتبری ندارد با \acro{IMSI}، معرفی می‌کند.
  \item اگر \acro{VLR} سه‌تایی استفاده‌نشده‌ای نداشته باشد، با پیام \acro{MAP}
        \msg{Send Authentication Info} دسته‌ای از سه‌تایی‌ها را از \acro{HLR/AuC} درخواست
        می‌کند و \acro{AuC} چند سه‌تایی بازمی‌گرداند.
  \item \acro{MSC/VLR} پیام \msg{Authentication Request} را تنها با \lr{RAND} می‌فرستد.
  \item \acro{SIM} مقادیر \lr{SRES}$'$ و $K_c$ را از \cref{eq:a3a8} محاسبه می‌کند و \acro{MS}
        مقدار \lr{SRES}$'$ را بازمی‌گرداند.
  \item \acro{VLR} مقدار \lr{SRES}$'$ را با \lr{SRES} مقایسه می‌کند. در صورت برابری، با فرستادن
        $K_c$ به \acro{BSS} در پیام \lr{BSSMAP} \msg{Cipher Mode Command} دستور رمزنگاری می‌دهد.
  \item \acro{BTS} پیام \msg{Ciphering Mode Command} را می‌فرستد؛ \acro{MS} رمز \acro{A5} را با
        $K_c$ آغاز می‌کند و با \msg{Ciphering Mode Complete} پاسخ می‌دهد. از این پس \acro{Um}
        رمزنگاری‌شده است.
  \item \acro{VLR} ممکن است یک \acro{TMSI} جدید تخصیص دهد که به‌صورت رمزشده ارسال می‌شود تا
        دفعهٔ بعد \acro{IMSI} محافظت شود.
\end{steps}
\end{procedurebox}

\begin{acnfigure}
  \resizebox{\linewidth}{!}{\acnfig{lecture-01/gsm-authentication}}
  \caption{توالی پیام‌های احراز هویت و رمزنگاری \acro{GSM} (\cref{proc:auth}).}
  \label{fig:auth}
\end{acnfigure}

\begin{procedurebox}{اتصال \acro{GPRS} و فعال‌سازی زمینهٔ \acro{PDP}}\label{proc:pdp}
\begin{steps}
  \item \textbf{اتصال \acro{GPRS}.} \acro{MS} پیام \msg{Attach Request} را به \acro{SGSN}
        می‌فرستد. \acro{SGSN} هویت \acro{MS} را احراز می‌کند (سه‌تایی‌ها از \acro{HLR} و از
        طریق \acro{Gr})، با \msg{Update GPRS Location} خود را در \acro{HLR} ثبت می‌کند، اشتراک
        را با \msg{Insert Subscriber Data} دریافت می‌کند و \msg{Attach Accept} را همراه با یک
        \acro{P-TMSI} بازمی‌گرداند. اکنون \acro{MS} در حالت \lr{READY} است اما نشانی \acro{IP} ندارد.
  \item \textbf{درخواست نشست.} \acro{MS} پیام \msg{Activate PDP Context Request} را همراه با
        \acro{NSAPI} (شناسهٔ زمینه)، نام نقطهٔ دسترسی (\acro{APN}) درخواستی و کیفیت خدمات
        می‌فرستد.
  \item \textbf{انتخاب دروازه.} \acro{SGSN} اشتراک را بررسی و \acro{APN} را با \acro{DNS} به
        نشانی یک \acro{GGSN} تبدیل می‌کند.
  \item \textbf{برقراری تونل.} \acro{SGSN} پیام \lr{GTP-C} \msg{Create PDP Context Request} را
        می‌فرستد؛ \acro{GGSN} یک نشانی \acro{IP} تخصیص می‌دهد (یا آن را از یک کارگزار
        \lr{RADIUS/DHCP} می‌گیرد) و شناسهٔ نقطهٔ پایانی تونل خود (\acro{TEID}) را بازمی‌گرداند.
  \item \textbf{تکمیل.} \acro{SGSN} پیام \msg{Activate PDP Context Accept} را همراه با نشانی
        \acro{IP} می‌فرستد. اکنون بسته‌های کاربر جریان می‌یابند: میان \acro{MS} و \acro{SGSN} روی
        \lr{LLC/SNDCP} و میان \acro{SGSN} و \acro{GGSN} درون یک تونل \lr{GTP-U} که با \acro{TEID}ها
        شناسایی می‌شود.
\end{steps}
\end{procedurebox}

\begin{acnfigure}
  \resizebox{\linewidth}{!}{\acnfig{lecture-01/pdp-activation}}
  \caption{اتصال \acro{GPRS} و سپس فعال‌سازی زمینهٔ \acro{PDP} (\cref{proc:pdp}).
  آبی خط‌چین: سیگنالینگ؛ فیروزه‌ای توپر: صفحهٔ کاربر.}
  \label{fig:pdp}
\end{acnfigure}

این رویه اصلی را نشان می‌دهد که دوباره در \acro{LTE} (حامل پیش‌فرض، درس ۹) و \lr{5G}
(نشست \acro{PDU}، درس ۲۰) با آن روبه‌رو خواهیم شد: \emph{متصل بودن} (شناخته‌شده برای شبکه)
و \emph{داشتن اتصال‌پذیری} (نشانی \acro{IP} و تونل) دو حالت جدا هستند که با دو رویهٔ جدا
برقرار می‌شوند. \acro{LTE} بعدها برای کاهش تأخیر این دو را برای نخستین اتصال \acro{PDN}
در هم ادغام کرد.

% ---------------------------------------------------------------------
\section{مثال‌های حل‌شده}\label{sec:examples}

\begin{examplebox}{ظرفیت یک اپراتور \lr{GSM-900}}\label{ex:capacity}
\textbf{مسئله.} اپراتوری \qty{10}{\mega\hertz} طیف جفت‌شده در \lr{GSM-900} دارد و با اندازهٔ
خوشهٔ $N=4$ طراحی می‌کند. هر سلول چند کانال صوتی تمام‌نرخ دارد؟

\textbf{حل.} تعداد حامل‌ها: $\floor{\qty{10}{\mega\hertz}/\qty{200}{\kilo\hertz}} = 50$؛ یکی به‌عنوان
محافظ کنار گذاشته می‌شود، پس ۴۹ حامل قابل استفاده است. برای هر سلول: $\floor{49/4} = 12$
حامل، یعنی $12\times 8 = 96$ شیار زمانی. با رزرو یک شیار برای کانال پخش و کنترل مشترک و
یکی برای سیگنالینگ \acro{SDCCH}، حدود ۹۴ کانال ترافیک باقی می‌ماند.

\textbf{تفسیر.} ۹۴ تماس هم‌زمان در هر سلول حد سخت سوئیچ مداری است. اینکه این کانال‌ها در
یک احتمال انسداد معین چقدر \emph{ترافیک} (بر حسب ارلانگ) حمل می‌کنند، در درس ۲ با فرمول
\lr{Erlang-B} محاسبه می‌شود.
\end{examplebox}

\begin{examplebox}{پیش‌فرستی زمانی}\label{ex:ta}
\textbf{مسئله.} \acro{BTS} اندازه می‌گیرد که بُرست‌های یک \acro{MS} به اندازهٔ
\qty{55}{\micro\second} دیر می‌رسند. چه مقدار \acro{TA} اعلام می‌شود و \acro{MS} چقدر دور است؟

\textbf{حل.} رفت‌وبرگشت \qty{55}{\micro\second} طول می‌کشد؛ پس
$\mathrm{TA} = \mathrm{round}(55/3.692) = 15$ و
$d \approx 15\times\qty{553}{\metre} \approx \qty{8.3}{\kilo\metre}$.
(به‌طور مستقیم: $d = c\,\tau/2 = 3\times10^8\times55\times10^{-6}/2 = \qty{8.25}{\kilo\metre}$.)

\textbf{تفسیر.} خطای کوانتش \acro{TA} تا نصف یک گام، حدود \qty{277}{\metre}، است؛ بنابراین
\acro{TA} تخمینی تقریبی از فاصلهٔ کاربر نیز به شبکه می‌دهد که پایهٔ مکان‌یابی اولیهٔ
\lr{cell-ID+TA} است.
\end{examplebox}

\begin{examplebox}{انتخاب طرح کدگذاری \acro{GPRS}}\label{ex:cs}
\textbf{مسئله.} در یک $C/I$ معین، کاربری با \lr{CS-2} مقدار $p_{\mathrm{BLER}} = 0.05$ و با
\lr{CS-4} مقدار $p_{\mathrm{BLER}} = 0.45$ را مشاهده می‌کند. با چهار شیار زمانی، بر اساس
\cref{eq:gprsrate} کدام طرح گذردهی بیشتری دارد؟

\textbf{حل.}
\begin{align*}
  R_{\mathrm{CS\text{-}2}} &= 4\times13.4\times(1-0.05) = \qty{50.9}{\kilo\bit\per\second},\\
  R_{\mathrm{CS\text{-}4}} &= 4\times21.4\times(1-0.45) = \qty{47.1}{\kilo\bit\per\second}.
\end{align*}
\textbf{تفسیر.} \lr{CS-2} با وجود نرخ اسمی کمتر برنده است. تطبیق پیوند دقیقاً انتخاب
طرحی است که این حاصل‌ضرب را بیشینه می‌کند؛ همین منطق با جدول‌های \acro{CQI} و \acro{MCS}
در \acro{LTE} (درس ۷) دوباره ظاهر می‌شود.
\end{examplebox}

% ---------------------------------------------------------------------
\section{مقایسهٔ فناوری‌ها}\label{sec:comparison}
\begin{table}[htbp]
\centering\small
\caption{مقایسهٔ \acro{GSM}، \acro{GPRS} و \acro{EDGE}. نرخ‌ها بیشینهٔ نرخ \acro{RLC} برای هر
شیار زمانی یا برای هر کاربر با هشت شیار هستند؛ نرخ‌های قابل دستیابی کمترند.}\label{tab:compare}
\rowcolors{2}{acnmist}{white}
\begin{tabularx}{\linewidth}{@{}l Y Y Y@{}}
\toprule
\rowcolor{white} & \acro{GSM} (دادهٔ مداری) & \acro{GPRS} & \acro{EDGE} (\acro{EGPRS}) \\
\midrule
نوع سوئیچ          & مداری                 & بسته‌ای                  & بسته‌ای \\
مدولاسیون          & \acro{GMSK}            & \acro{GMSK}              & \acro{GMSK} و \lr{8-PSK} \\
بیت در هر نماد     & ۱                      & ۱                        & ۱ یا ۳ \\
نرخ هر شیار        & \lr{9.6} یا \lr{14.4\,kbit/s} & \lr{9.05--21.4\,kbit/s} & \lr{8.8--59.2\,kbit/s} \\
بیشینه (۸ شیار)    & --                     & \lr{171.2\,kbit/s}        & \lr{473.6\,kbit/s} \\
استفاده از منابع   & رزروشده برای تماس      & اشتراکی، بنا به تقاضا    & اشتراکی، بنا به تقاضا \\
ارسال مجدد         & ندارد (فقط \acro{FEC})  & \acro{ARQ} در \acro{RLC}  & \acro{ARQ} + افزونگی افزایشی \\
گره‌های هسته       & \lr{MSC, GMSC}         & \lr{SGSN, GGSN}           & \lr{SGSN, GGSN} (بدون تغییر) \\
مبنای صورت‌حساب    & زمان                   & حجم                       & حجم \\
\bottomrule
\end{tabularx}
\end{table}

\paragraph{مصالحه‌ها} \acro{GPRS} داده را ارزان کرد (صورت‌حساب بر اساس حجم، بدون مدار بیکار)،
اما تأخیر صف و گذردهی متغیر را به همراه آورد؛ زمان رفت‌وبرگشت چندصد میلی‌ثانیه‌ای رایج بود.
\acro{EDGE} با یک ارتقای نرم‌افزاری و فرستنده‌گیرنده بیشینهٔ نرخ را سه برابر کرد، اما
\lr{8-PSK} پوش ثابت ندارد و به تقویت‌کننده‌های توان خطی‌تر نیاز دارد و طرح‌های بالای آن تنها
نزدیک \acro{BTS} قابل استفاده‌اند. هیچ‌کدام محدودیت بنیادی را تغییر نداد: داده بر حامل‌های
\qty{200}{\kilo\hertz} سوار است که برای صدا طراحی شده‌اند. شکستن این محدودیت به واسط رادیویی
جدیدی نیاز داشت: \acro{WCDMA} در نسل سوم (درس ۳).

% ---------------------------------------------------------------------
\section{نکته‌های کلیدی}\label{sec:takeaways}
\begin{takeawaybox}
\begin{itemize}
  \item \acro{GSM} یک استاندارد \emph{سامانه‌ای} است: واسط‌های باز (\acro{A}، \acro{Um}) و هویت
        جداشدنی (\lr{SIM/IMSI}) شبکه‌های چندسازنده‌ای و رومینگ را ممکن کردند.
  \item همهٔ زمان‌بندی \acro{GSM} از $R = \qty{270.833}{\kilo\bit\per\second}$ نتیجه می‌شود: بیت
        \qty{3.69}{\micro\second}، شیار \qty{577}{\micro\second}، قاب \qty{4.615}{\milli\second}،
        چندقاب ۲۶تایی \qty{120}{\milli\second}؛ یک \acro{TCH} تمام‌نرخ \qty{22.8}{\kilo\bit\per\second}
        ناخالص حمل می‌کند.
  \item پیش‌فرستی زمانی ۶ بیتی شعاع یک سلول عادی را به حدود \qty{35}{\kilo\metre} محدود می‌کند.
  \item احراز هویت از سه‌تایی‌های $(\text{RAND},\text{SRES},K_c)$ استفاده می‌کند؛ یک‌طرفه است و
        رمزنگاری در \acro{BTS} پایان می‌یابد.
  \item \acro{GPRS} گره‌های \acro{PCU}، \acro{SGSN} و \acro{GGSN} را می‌افزاید؛ زمینهٔ \acro{PDP} و
        تونل \acro{GTP} نیاکان حامل \acro{LTE} و نشست \acro{PDU} در \lr{5G} هستند.
  \item \acro{EDGE} نرخ بیت هر نماد را افزایش می‌دهد (\lr{8-PSK}) و افزونگی افزایشی را
        می‌افزاید که شکل اولیهٔ \acro{HARQ} است.
  \item مدل ردیف اول $C/I \approx (\sqrt{3N})^n/6$ بازاستفاده، افت مسیر و تداخل را به هم پیوند
        می‌دهد؛ این یک تخمین اولیهٔ طراحی است.
\end{itemize}
\end{takeawaybox}

% ---------------------------------------------------------------------
\section{تمرین‌ها}\label{sec:exercises}

\begin{exercisebox}{1}{معماری}\label{exr:arch}
توضیح دهید چرا \acro{HLR} نشانی \acro{VLR} را ذخیره می‌کند اما سلولی را که مشترک در آن
قرار دارد ذخیره نمی‌کند. پیامد سیگنالینگی ذخیرهٔ سلول چه می‌بود؟
\end{exercisebox}

\begin{exercisebox}{1}{امنیت}\label{exr:sec}
دو حمله را که احراز هویت یک‌طرفهٔ \acro{GSM} ممکن می‌سازد نام ببرید و بگویید کدام سازوکار
بعدی (درس ۱۱) جلوی هر یک را می‌گیرد.
\end{exercisebox}

\begin{exercisebox}{2}{کانال نیم‌نرخ}\label{exr:halfrate}
یک کانال ترافیک نیم‌نرخ در هر دو قاب ترافیک از یک شیار زمانی استفاده می‌کند. با استفاده
از \cref{eq:tchrate} نرخ بیت ناخالص آن را محاسبه کنید و توضیح دهید چگونه ظرفیت صوتی سلول
\cref{ex:capacity} را دو برابر می‌کند.
\end{exercisebox}

\begin{exercisebox}{2}{طراحی بازاستفاده}\label{exr:reuse}
با استفاده از \cref{eq:ci}، کوچک‌ترین اندازهٔ خوشهٔ معتبر $N$ را بیابید که
$C/I\ge\qty{12}{\dB}$ را برای (الف) $n=4$ و (ب) $n=3$ تضمین کند.
\end{exercisebox}

\begin{exercisebox}{2}{گذردهی \acro{EDGE}}\label{exr:edge}
کاربری پنج شیار زمانی دارد. با \lr{MCS-9} نسبت خطای بلوک 0.30 و با \lr{MCS-7}
(\qty{44.8}{\kilo\bit\per\second} برای هر شیار) برابر 0.08 است. تطبیق پیوند کدام طرح را
باید انتخاب کند و گذردهی مورد انتظار چقدر است؟
\end{exercisebox}

\begin{exercisebox}{3}{بودجهٔ سربار}\label{exr:overhead}
با شروع از \cref{eq:tb,eq:tframe}، محاسبه کنید چه کسری از نرخ بیت ناخالص یک شیار زمانی
صرف این موارد می‌شود: (۱) دورهٔ محافظ، (۲) بیت‌های انتهایی و پرچم‌های ربایش، (۳) دنبالهٔ
آموزشی، (۴) قاب‌های \acro{SACCH} و بیکار. مجموع را با سربار پیشوند چرخه‌ای \acro{LTE}
(حدود $7\%$) مقایسه کنید و توضیح دهید چرا این مقایسه کاملاً منصفانه نیست.
\end{exercisebox}

\begin{exercisebox}{3}{تعمیم مدل \lr{C/I}}\label{exr:ci3}
\cref{eq:ci} را برای سلولی که به سه بخش \ang{120} تقسیم شده تعمیم دهید، به‌طوری که تنها دو
تداخل‌کنندهٔ ردیف اول در لوب اصلی آنتن بخش قرار گیرند. عبارت جدید و بهره بر حسب دسی‌بل را
استخراج کنید و بحث کنید نتیجه کدام فرض‌های مدل اصلی را به ارث می‌برد.
\end{exercisebox}

\subsection*{پاسخ‌ها}
\addcontentsline{toc}{subsection}{پاسخ‌ها}

\begin{solutionbox}{\ref{exr:arch}}
سلول کاربر متحرک هر چند دقیقه یک بار عوض می‌شود؛ اگر \acro{HLR} آن را دنبال می‌کرد، هر تغییر
سلول هر مشترک در جهان سیگنالینگی به شبکهٔ خانگی، اغلب به‌صورت بین‌المللی، تولید می‌کرد.
ذخیرهٔ تنها نشانی \acro{VLR} به‌روزرسانی‌های \acro{HLR} را به تغییر ناحیهٔ \acro{MSC} محدود
می‌کند و ناحیهٔ مکانی ذخیره‌شده در \acro{VLR} فراخوانی رادیویی را به گروهی از سلول‌ها محدود
می‌کند. این سلسله‌مراتب دوسطحی (\lr{HLR $\to$ VLR $\to$ LA}) هستهٔ مدیریت تحرک \acro{GSM} است.
\end{solutionbox}

\begin{solutionbox}{\ref{exr:sec}}
(۱) یک ایستگاه پایهٔ جعلی می‌تواند \acro{IMSI} را به‌صورت آشکار درخواست کند و کاربر را
ردیابی کند (\lr{IMSI catching})؛ \lr{5G} هویت دائمی را با \acro{SUCI} پنهان می‌کند و
\lr{LTE/UMTS} این افشا را کاهش می‌دهند. (۲) ایستگاه پایهٔ جعلی می‌تواند \acro{MS} را بدون
احراز هویت بپذیرد، دستور «بدون رمزنگاری» یا رمز ضعیف بدهد و ترافیک را رله کند (حملهٔ
مرد میانی). احراز هویت دوطرفهٔ \lr{UMTS/LTE} با نشانهٔ \acro{AUTN} (درس ۱۱) به \acro{USIM}
امکان می‌دهد بررسی کند که چالش را شبکهٔ خانگی خودش تولید کرده است، و حفاظت از یکپارچگی
سیگنالینگ از دست‌کاری فرمان حالت امنیتی جلوگیری می‌کند.
\end{solutionbox}

\begin{solutionbox}{\ref{exr:halfrate}}
کانال نیم‌نرخ ۱۲ قاب از ۲۴ قاب ترافیک هر چندقاب ۲۶تایی را می‌گیرد؛ پس
$R_{\mathrm{TCH/H}} = 114\times 12/0.120 = \qty{11.4}{\kilo\bit\per\second}$. دو کاربر نیم‌نرخ
یک شیار زمانی را به اشتراک دارند؛ بنابراین ۹۴ شیار \cref{ex:capacity} می‌توانند تا ۱۸۸ تماس
را حمل کنند، به بهای کدک گفتار با نرخ (و کیفیت) کمتر.
\end{solutionbox}

\begin{solutionbox}{\ref{exr:reuse}}
باید $(\sqrt{3N})^n/6 \ge 10^{1.2} = 15.85$ باشد، یعنی $(3N)^{n/2}\ge 95.1$.
(الف) $n=4$: $(3N)^2\ge 95.1 \Rightarrow N\ge 3.25$؛ کوچک‌ترین $N$ معتبر ۴ است
($C/I=\qty{13.8}{\dB}$). (ب) $n=3$: $(3N)^{1.5}\ge95.1 \Rightarrow 3N\ge 20.8 \Rightarrow N\ge 6.94$؛
کوچک‌ترین $N$ معتبر ۷ است ($C/I = 10\log_{10}(21^{1.5}/6) = \qty{12.05}{\dB}$).
اندازه‌های معتبر عبارت‌اند از $N = i^2+ij+j^2 \in\{1,3,4,7,9,12,\dots\}$.
\end{solutionbox}

\begin{solutionbox}{\ref{exr:edge}}
$R_{\mathrm{MCS\text{-}9}} = 5\times59.2\times0.70 = \qty{207.2}{\kilo\bit\per\second}$ و
$R_{\mathrm{MCS\text{-}7}} = 5\times44.8\times0.92 = \qty{206.1}{\kilo\bit\per\second}$.
این دو عملاً برابرند؛ \lr{MCS-9} از نظر گذردهی اندکی بهتر است، اما \lr{MCS-7} ارسال مجدد
کمتر و در نتیجه تأخیر و لرزش کمتری دارد. با افزونگی افزایشی، گذردهی مؤثر \lr{MCS-9} بیشتر
از پیش‌بینی مدل ساده خواهد بود، زیرا بلوک‌های ناموفق نیز اطلاعات نرم فراهم می‌کنند.
\end{solutionbox}

\begin{solutionbox}{\ref{exr:overhead}}
در هر شیار ۱۵۶٫۲۵ بیتی: محافظ $8.25/156.25 = 5.3\%$؛ بیت‌های انتهایی و پرچم‌ها
$(6+2)/156.25 = 5.1\%$؛ دنبالهٔ آموزشی $26/156.25 = 16.6\%$؛ داده $114/156.25 = 73.0\%$. از
۲۶ قاب، ۲ قاب ترافیک حمل نمی‌کنند؛ پس کسر مفید $0.730\times24/26 = 67.4\%$ است:
$270.833/8\times0.674 = \qty{22.8}{\kilo\bit\per\second}$ که با \cref{eq:tchrate} سازگار است.
سربار کل حدود $33\%$ است. پیشوند چرخه‌ای \acro{LTE} حدود $7\%$ هزینه دارد، اما \acro{LTE}
منابعی را نیز صرف سیگنال‌های مرجع، ناحیهٔ کنترل و باندهای محافظ می‌کند؛ یک مقایسهٔ
منصفانه باید همهٔ آن‌ها را بشمارد (درس ۷).
\end{solutionbox}

\begin{solutionbox}{\ref{exr:ci3}}
با دو تداخل‌کننده، $C/I \approx (\sqrt{3N})^n/2$ که سه برابر (\qty{4.8}{\dB}) بهتر از
\cref{eq:ci} است. برای $N=4$ و $n=4$: $10\log_{10}(144/2) = \qty{18.6}{\dB}$. این نتیجه همچنان
شش‌ضلعی‌های ایده‌آل، توان‌های برابر، نبود سایه و کاربر بدترین‌حالت در لبهٔ سلول را فرض
می‌کند؛ افزون بر آن الگوهای آنتن ایده‌آل بدون نشت لوب پشتی را فرض می‌کند. بخش‌بندی، بازده
ترانکینگ (هر بخش کانال‌های کمتری دارد) را با کاهش تداخل معاوضه می‌کند.
\end{solutionbox}

% ---------------------------------------------------------------------
\section{مطالعهٔ بیشتر}\label{sec:reading}
معماری همهٔ سامانه‌های \acro{3GPP} در \cite{ts23002} خلاصه شده است؛ ساختار رادیویی \acro{GSM}
در \cite{ts45002} و کارکردهای امنیتی در \cite{ts43020} مشخص شده‌اند. رویه‌های مرحلهٔ ۲
\acro{GPRS} در \cite{ts23060} و پیام‌های لایهٔ ۳ در \cite{ts24008} آمده‌اند. کتاب
\lr{Eberspächer} و همکاران \cite{eberspaecher2009gsm} همچنان کامل‌ترین توصیف یک‌جای \acro{GSM}
و \acro{GPRS} است و \lr{Rappaport} \cite{rappaport2002wireless} مدل بازاستفادهٔ سلولی
به‌کاررفته در \cref{sec:reuse} را استخراج می‌کند. \acro{MAP} در \cite{ts29002} و هویت‌ها در
\cite{ts23003} مشخص شده‌اند.

\nocite{shannon1948mathematical}
\begin{latin}
\printbibliography[heading=subbibliography,title={References}]
\end{latin}

\end{document}
